% Generated by report/tools/tables.py from the result files of the repository. Do not edit.
\begin{table}[tbp]
\centering
\footnotesize
\setlength{\tabcolsep}{3pt}
\caption[Grille de Pineau : 2018 contre RL classique 2026]{Grille de Pineau (critères Q2 à Q8) et critères X tirés de sa conférence ICLR 2018, lus avec la même grille et les mêmes lecteurs sur deux échantillons : notre lecture de 50 articles de 2018 (ICML, NeurIPS, ICLR) et 50 articles de RL classique de 2026 (14 du premier tirage, 36 de la file A3). C'est la comparaison principale, la seule qui fonde un verdict (écart 20 de \file{study/PROTOCOL.md}). $k/n$ : articles qui remplissent le critère sur les articles où il s'applique (articles à expériences, Q1 différent de « non ») ; la part en pourcentage porte son intervalle de Wilson à 95~\%. Différence : RL classique 2026 moins 2018, en points de pourcentage, avec l'intervalle de Newcombe à 95~\% ; $p$ : test exact de Fisher bilatéral ; $p$ Holm : correction de Holm sur les onze critères de la grille v1. \textbf{En gras}, les critères dont la différence passe la correction de Holm ($p$ Holm $<$ 0,05). Pour X3, une part plus basse est un progrès ; Q2 (réseaux de neurones) décrit la population étudiée, pas une pratique de reproductibilité. Les chiffres de Pineau elle-même sont au tableau~\ref{tab:sondage-pineau-contexte}. Source : \file{study/extension\_stats.json} (script \file{study/extension\_stats.py}, écrit avant toute lecture des articles ajoutés).}
\label{tab:sondage-pineau}
\begin{tabularx}{\linewidth}{@{}l>{\raggedright\arraybackslash}Xr r@{\hspace{0.28em}}l r r@{\hspace{0.28em}}l r@{\hspace{0.28em}}l rr@{}}
\toprule
 & & \multicolumn{3}{c}{\makecell{Notre lecture\\{}de 2018}} & \multicolumn{3}{c}{\makecell{RL classique\\{}2026}} & \multicolumn{2}{c}{Différence} & &  \\
\cmidrule(lr){3-5}\cmidrule(lr){6-8}
Code & Critère & $k/n$ & \multicolumn{2}{c}{\% [Wilson]} & $k/n$ & \multicolumn{2}{c}{\% [Wilson]} & \multicolumn{2}{c}{[Newcombe]} & $p$ & $p$ Holm \\
\midrule
\textbf{Q2} & \textbf{Réseaux de neurones employés} & \textbf{38/48} & \textbf{79} & \textbf{[66\,;\,88]} & \textbf{49/50} & \textbf{98} & \textbf{[90\,;\,100]} & \textbf{+19} & \textbf{[+6\,;\,+32]} & \textbf{0,003} & \textbf{0,03} \\
Q3 & Tous les hyperparamètres de la méthode & 26/47 & 55 & [41\,;\,69] & 36/50 & 72 & [58\,;\,83] & +17 & [−2\,;\,+34] & 0,10 & 0,58 \\
Q4 & Tous les hyperparamètres des références & 17/48 & 35 & [23\,;\,50] & 14/49 & 29 & [18\,;\,42] & −7 & [−25\,;\,+11] & 0,52 & 1 \\
\textbf{Q5} & \textbf{Code en lien} & \textbf{15/48} & \textbf{31} & \textbf{[20\,;\,45]} & \textbf{37/50} & \textbf{74} & \textbf{[60\,;\,84]} & \textbf{+43} & \textbf{[+23\,;\,+58]} & \textbf{4\,×\,10\textsuperscript{−5}} & \textbf{0,0004} \\
Q6 & Choix des hyperparamètres décrit & 12/47 & 26 & [15\,;\,40] & 24/50 & 48 & [35\,;\,61] & +22 & [+3\,;\,+39] & 0,03 & 0,24 \\
Q7 & Évaluation sur conditions tenues à part & 16/48 & 33 & [22\,;\,47] & 22/49 & 45 & [32\,;\,59] & +12 & [−8\,;\,+30] & 0,30 & 1 \\
Q8 & Test statistique appliqué & 3/48 & 6 & [2\,;\,17] & 6/50 & 12 & [6\,;\,24] & +6 & [−7\,;\,+18] & 0,49 & 1 \\
\textbf{X2} & \textbf{Incertitude rapportée} & \textbf{30/48} & \textbf{63} & \textbf{[48\,;\,75]} & \textbf{44/50} & \textbf{88} & \textbf{[76\,;\,94]} & \textbf{+26} & \textbf{[+8\,;\,+41]} & \textbf{0,005} & \textbf{0,04} \\
X3 & Seuls les meilleurs essais gardés & 9/48 & 19 & [10\,;\,32] & 7/50 & 14 & [7\,;\,26] & −5 & [−20\,;\,+10] & 0,59 & 1 \\
X4 & Références réglées avec le même soin & 6/48 & 13 & [6\,;\,25] & 8/49 & 16 & [9\,;\,29] & +4 & [−11\,;\,+18] & 0,77 & 1 \\
\textbf{X7} & \textbf{Calcul chiffré} & \textbf{13/48} & \textbf{27} & \textbf{[17\,;\,41]} & \textbf{29/50} & \textbf{58} & \textbf{[44\,;\,71]} & \textbf{+31} & \textbf{[+11\,;\,+47]} & \textbf{0,002} & \textbf{0,02} \\
\bottomrule
\end{tabularx}
\end{table}
